\chapter{2007年广东卷（文）}

\section{单选题}

\begin{problem}
已知集合 \(M = \{x \mid 1 + x > 0\}\)，\(N = \left\{x \mid \frac{1}{1 - x} > 0\right\}\)，则 \(M \cap N =\)（\quad）

\choices
    {\(\{x\mid -1\leqslant x < 1\}\)}
    {\( \{x \mid x > 1\} \)}
    {\(\{x\mid -1 < x < 1\}\)}
    {\( \{x \mid x \geqslant -1\} \)}
\end{problem}

\begin{problem}
若复数 \((1 + b\i)(2 + \i)\) 是纯虚数（\(\i\) 是虚数单位，\(b\) 是实数），则 \(b =\)（\quad）

\choices
    {\(-2\)}
    {\( -\frac{1}{2} \)}
    {\( \frac{1}{2} \)}
    {\(2\)}
\end{problem}

\begin{problem}
若函数 \( f(x) = x^3 \)（\(x \in \R\)），则函数 \( y = f(-x) \) 在其定义域上是（\quad）

\choices
    {单调递减的偶函数}
    {单调递减的奇函数}
    {单调递增的偶函数}
    {单调递增的奇函数}
\end{problem}

\begin{problem}
若向量 \(\bs{a}\)，\(\bs{b}\) 满足 \(|\bs{a}| = |\bs{b}| = 1\)，\(\bs{a}\) 与 \(\bs{b}\) 的夹角为 \(60^\circ\)，则 \(\bs{a} \cdot \bs{a} + \bs{a} \cdot \bs{b} =\)（\quad）

\choices
    {\(\frac{1}{2}\)}
    {\(\frac{3}{2}\)}
    {\(1 + \frac{\sqrt{3}}{2}\)}
    {\(2\)}
\end{problem}

\begin{problem}
客车从甲地以 \(60\,\mathrm{km/h}\) 的速度匀速行驶 1 小时到达乙地，在乙地停留了半小时，然后以 \(80\,\mathrm{km/h}\) 的速度匀速行驶 1 小时到达丙地. 下列描述客车从甲地出发，经过乙地，最后到达丙地所经过的路程 \(s\) 与时间 \(t\) 之间关系的图象中，正确的是（\quad）

\bitmapfigure[max width=0.88\linewidth]{img/2007/guangdong_liberal/q5_options.png}
\end{problem}

\begin{problem}
若 \(l\)，\(m\)，\(n\) 是互不相同的空间直线，\(\alpha\)，\(\beta\) 是不重合的平面，则下列命题中为真命题的是（\quad）

\choices
    {若 \(\alpha \parallel \beta, l \subset \alpha, n \subset \beta\)，则 \(l \parallel n\)}
    {若 \(\alpha \perp \beta\)，\(l \subset \alpha\)，则 \(l \perp \beta\)}
    {若 \(l \perp n\)，\(m \perp n\)，则 \(l \parallel m\)}
    {若 \(l \perp \alpha\)，\(l \parallel \beta\)，则 \(\alpha \perp \beta\)}
\end{problem}

\begin{problem}
图1是某县参加2007年高考的学生身高条形统计图，从左到右的条形表示的学生人数依次记为 \(A_{1}\)，\(A_{2}\)，\(\cdots\)，\(A_{10}\)（如 \(A_{2}\) 表示身高（单位：cm）在 \([150,155)\) 内的学生人数）. 图2是统计图1中身高在一定范围内学生人数的一个算法流程图. 现要统计身高在 \(160 \sim 180 \mathrm{~cm}\)（含 \(160 \mathrm{~cm}\)，不含 \(180 \mathrm{~cm}\)）的学生人数，那么在流程图中的判断框内应填写的条件是（\quad）

\bitmapfigure[max width=0.9\linewidth]{img/2007/guangdong_liberal/q7_fig1.png}

\texfigure[max width=0.9\linewidth]{img_repaint/2007/guangdong_liberal/q7_fig2.tex}

\choices
    {\(i < 9\)}
    {\(i < 8\)}
    {\(i < 7\)}
    {\(i < 6\)}
\end{problem}

\begin{problem}
在一个袋子中装有分别标注数字 1， 2， 3， 4， 5 的五个小球，这些小球除标注的数字外完全相同. 现从中随机取出 2 个小球，则取出的小球标注的数字之和为 3 或 6 的概率是（\quad）

\choices
    {\(\frac{3}{10}\)}
    {\(\frac{1}{5}\)}
    {\(\frac{1}{10}\)}
    {\(\frac{1}{12}\)}
\end{problem}

\begin{problem}
已知简谐运动 \( f(x) = 2\sin \left(\frac{\pi}{3} x + \varphi\right) \left(|\varphi| < \frac{\pi}{2}\right) \) 的图象经过点 \((0,1)\)，则该简谐运动的最小正周期 \(T\) 和初相 \(\varphi\) 分别为（\quad）

\choices
    {\(T = 6, \varphi = \frac{\pi}{6}\)}
    {\(T = 6\)，\(\varphi = \frac{\pi}{3}\)}
    {\(T = 6\pi\)，\(\varphi = \frac{\pi}{6}\)}
    {\(T = 6\pi\)，\(\varphi = \frac{\pi}{3}\)}
\end{problem}

\begin{problem}
如图是某汽车维修公司的维修点环形分布图. 公司在年初分配给 \(A\)，\(B\)，\(C\)，\(D\) 四个维修点某种配件各 50 件. 在使用前发现需将 \(A\)，\(B\)，\(C\)，\(D\) 四个维修点的这批配件分别调整为 40， 45， 54， 61 件，但调整只能在相邻维修点之间进行，那么要完成上述调整，最少的调动件次（\(n\) 件配件从一个维修点调整到相邻维修点的调动件次为 \(n\)）为（\quad）

\bitmapfigure[max width=0.38\linewidth]{img/2007/guangdong_liberal/q10_fig1.png}

\choices
    {18}
    {17}
    {16}
    {15}
\end{problem}

\section{填空题}

\begin{problem}
在平面直角坐标系 \(xOy\) 中，已知抛物线关于 \(x\) 轴对称，顶点在原点 \(O\)，且过点 \(P(2,4)\)，则该抛物线的方程是\fillinblank{}.
\end{problem}

\begin{problem}
函数 \( f(x) = x \ln x \)（\(x > 0\)）的单调递增区间是\fillinblank{}.
\end{problem}

\begin{problem}
已知数列 \(\{a_{n}\}\) 的前 \(n\) 项和 \(S_{n}=n^{2}-9n\)，则其通项 \(a_{n}=\)\fillinblank{}；若它的第 \(k\) 项满足 \(5<a_{k}<8\)，则 \(k=\)\fillinblank{}.
\end{problem}

\begin{problem}
在极坐标系中，直线 \(l\) 的方程为 \(\rho \sin \theta = 3\)，则点 \(\left(2, \frac{\pi}{6}\right)\) 到直线 \(l\) 的距离为\fillinblank{}.
\end{problem}

\begin{problem}
如图所示，圆 \(O\) 的直径 \(AB = 6\)，\(C\) 为圆周上一点，\(BC = 3\).过 \(C\) 作圆的切线 \(l\)，过 \(A\) 作 \(l\) 的垂线 \(AD\)，垂足为 \(D\)，则 \(\angle DAC =\)\fillinblank{}.

\bitmapfigure[max width=0.38\linewidth]{img/2007/guangdong_liberal/q15_fig1.png}
\end{problem}

\section{解答题}

\begin{problem}
已知 \(\triangle ABC\) 顶点的直角坐标分别为 \(A(3,4)\)，\(B(0,0)\)，\(C(c,0)\).

\begin{enumerate}
\item 若 \(\overrightarrow{AB} \cdot \overrightarrow{AC} = 0\)，求 \(c\) 的值；
\item 若 \(c = 5\)，求 \(\sin \angle A\) 的值.
\end{enumerate}
\end{problem}

\begin{problem}
已知某几何体的俯视图是如图所示的矩形，正视图（或称主视图）是一个底边长为 \(8\)，高为 \(4\) 的等腰三角形，侧视图（或称左视图）是一个底边长为 \(6\)，高为 \(4\) 的等腰三角形.

\begin{enumerate}
\item 求该几何体的体积 \(V\)；
\item 求该几何体的侧面积 \(S\).
\end{enumerate}

\bitmapfigure[max width=0.35\linewidth]{img/2007/guangdong_liberal/q17_fig1.png}

\end{problem}

\begin{problem}
下表提供了某厂节能降耗技术改造后生产甲产品过程中记录的产量 \(x\)（吨）与相应的生产能耗 \(y\)（吨标准煤）的几组对照数据.

\begin{center}
\begin{tabular}{c|cccc}
\(x\) & 3 & 4 & 5 & 6 \\
\hline
\(y\) & 2.5 & 3 & 4 & 4.5
\end{tabular}
\end{center}

\begin{enumerate}
\item 请画出上表的散点图；
\item 请根据上表提供的数据，用最小二乘法求出 \(y\) 关于 \(x\) 的线性回归方程 \(y = \hat{b}x + \hat{a}\)；
\item 已知该厂技改前 100 吨甲产品的生产能耗为 90 吨标准煤. 试根据前一问求出的线性回归方程，预测生产 100 吨甲产品的生产能耗比技改前降低多少吨标准煤？（参考数值：\(3 \times 2.5 + 4 \times 3 + 5 \times 4 + 6 \times 4.5 = 66.5\)）
\end{enumerate}
\end{problem}

\begin{problem}
在平面直角坐标系 \(xOy\) 中，已知圆心在第二象限，半径为 \(2\sqrt{2}\) 的圆 \(C\) 与直线 \(y = x\) 相切于坐标原点 \(O\). 椭圆 \(\frac{x^2}{a^2} + \frac{y^2}{9} = 1\) 与圆 \(C\) 的一个交点到椭圆两焦点的距离之和为 \(10\).

\begin{enumerate}
\item 求圆 \(C\) 的方程；
\item 试探求 \(C\) 上是否存在异于原点的点 \(Q\)，使 \(Q\) 到椭圆右焦点 \(F\) 的距离等于线段 \(OF\) 的长. 若存在，请求出点 \(Q\) 的坐标；若不存在，请说明理由.
\end{enumerate}
\end{problem}

\begin{problem}
已知函数 \( f(x) = x^2 + x - 1 \)，\(\alpha\)，\(\beta\) 是方程 \(f(x) = 0\) 的两个根（\(\alpha > \beta\)），\(f'(x)\) 是 \(f(x)\) 的导数，设 \(a_{1} = 1\)，\(a_{n + 1} = a_{n} - \frac{f(a_{n})}{f'(a_{n})}\)（\(n = 1\)，\(2\)，\(\cdots\)）.

\begin{enumerate}
\item 求 \(\alpha\)，\(\beta\) 的值；
\item 已知对任意的正整数 \(n\)，都有 \(a_{n} > \alpha\)，记 \(b_{n} = \ln \frac{a_{n} - \beta}{a_{n} - \alpha}\)（\(n = 1\)，\(2\)，\(\cdots\)），求数列 \(\{b_{n}\}\) 的前 \(n\) 项和 \(S_{n}\).
\end{enumerate}
\end{problem}

\begin{problem}
已知 \(a\) 是实数，函数 \(f(x) = 2ax^2 + 2x - 3 - a\). 如果函数 \(y = f(x)\) 在区间 \([-1, 1]\) 上有零点，求 \(a\) 的取值范围.
\end{problem}
